% © 2026 Ing. David Jaroš — CC BY-NC-ND 4.0
%
% This work is licensed under a Creative Commons Attribution-NonCommercial-NoDerivatives
% 4.0 International License (CC BY-NC-ND 4.0).
%
% B_base non-perturbative approaches D1–D3
% Track: CORE — α derivation — v60 baseline
% Date: 2026-03-08

\ifdefined\INCLUDEMODE
  % Being included in another document — skip preamble
\else
  \documentclass[11pt,a4paper]{article}
  \usepackage{amsmath,amssymb,amsthm}
  \usepackage{hyperref}

  \newtheorem{theorem}{Theorem}
  \newtheorem{lemma}{Lemma}
  \newtheorem{proposition}{Proposition}
  \newtheorem{conjecture}{Conjecture}
  \theoremstyle{definition}
  \newtheorem{gap}{Gap}
  \theoremstyle{remark}
  \newtheorem{remark}{Remark}

  \title{$B_{\mathrm{base}}$ — Non-perturbative Approaches D1--D3\\
    \large Unitarity Constraint, Dimensional Transmutation,\\
    \large and Cartan--Killing Metric on $\mathrm{Im}(\mathbb{H})$}
  \author{Ing.\ David Jaroš}
  \date{2026-03-08 (v60)}

  \begin{document}
  \maketitle
\fi

%--------------------------------------------------------------------
\section{Recap: Known Results and Dead Ends}
\label{sec:recap}
%--------------------------------------------------------------------

\noindent\textbf{No circular reasoning.}
$\alpha^{-1} = 137$ and $n^* = 137$ are \emph{never} used as inputs anywhere
in this document.
The only inputs are: $N_{\mathrm{eff}} = 12$ (proved from $\mathbb{C}\otimes\mathbb{H}$
algebra), $\dim_\mathbb{R}(\mathrm{Im}\,\mathbb{H}) = 3$ (algebraic fact), and
the UBT kinetic action $S[\Theta]$.

\subsection{Established baseline (v60)}

\begin{center}
\begin{tabular}{lll}
  \hline
  Quantity & Value & Status \\
  \hline
  $B_0 = 2\pi N_{\mathrm{eff}}/3$ & $25.133$ & \textbf{[PROVED]} (one-loop, flat space) \\
  $B_{\mathrm{base}} = N_{\mathrm{eff}}^{3/2}$ & $41.569$ & [CONJECTURED] \\
  $B_{\mathrm{base}}/B_0 = (3/2\pi)\sqrt{N_{\mathrm{eff}}}$ & $1.6540$ & \textbf{algebraic identity} \\
  \hline
\end{tabular}
\end{center}

The algebraic ratio is an exact identity:
\begin{equation}
  \frac{B_{\mathrm{base}}}{B_0}
    = \frac{N_{\mathrm{eff}}^{3/2}}{2\pi N_{\mathrm{eff}}/3}
    = \frac{3}{2\pi}\,\sqrt{N_{\mathrm{eff}}}
    = \frac{3\sqrt{12}}{2\pi}
    \approx 1.6540.
  \label{eq:ratio_identity}
\end{equation}

\subsection{Perturbative dead ends (approaches A1--C2)}

Five perturbative approaches have been exhausted and documented; they must not be
repeated:

\begin{center}
\begin{tabular}{lll}
  \hline
  Approach & Dead-end reason & File \\
  \hline
  A1: Kaluza--Klein sum & Sum diverges; no natural cutoff & \texttt{b\_base\_spinor\_approach.tex} \\
  A2: Spectral zeta function & Exponent $\ne 3/2$ without free param & \texttt{b\_base\_hausdorff.tex} \\
  A3: Gauge orbit volume & Volume factors cancel; no $\sqrt{N}$ & \texttt{b\_base\_hausdorff.tex} \\
  C1: Seeley--DeWitt curvature & Correction $\propto \Lambda^{-2} \to 0$ & \texttt{b\_base\_delta\_d.tex} \\
  C2: Mode pair interference & Algebraic identity restatement & \texttt{b\_base\_delta\_d.tex} \\
  \hline
\end{tabular}
\end{center}

The present document investigates three \emph{non-perturbative} candidates
(D1--D3) that are algebraically independent of all previous attempts.

%--------------------------------------------------------------------
\section{Approach D1: Unitarity Constraint on $\mathrm{Im}(\mathbb{H})$}
\label{sec:D1_unitarity}
%--------------------------------------------------------------------

\textbf{Hypothesis.}
Not all $N_{\mathrm{eff}} = 12$ modes are independent after imposing the
unitarity condition $\Theta^\dagger\Theta = \mathbf{1}$ on the imaginary part
$\mathrm{Im}(\mathbb{H})$.
The reduced number of independent modes might equal $\sqrt{N_{\mathrm{eff}}} = 2\sqrt{3}$,
which would explain the factor in~\eqref{eq:ratio_identity}.

\subsection{Step 1: Unitarity condition for $\Theta \in \mathbb{C}\otimes\mathbb{H}$}

A general element of $\mathbb{C}\otimes\mathbb{H}$ can be written as
$\Theta = \theta_0 + \theta_1\,\mathbf{i} + \theta_2\,\mathbf{j} + \theta_3\,\mathbf{k}$
with $\theta_\mu \in \mathbb{C}$.
The unitarity condition $\Theta^\dagger\Theta = 1$ (where $\dagger$ denotes
quaternionic conjugation followed by complex conjugation) imposes
\begin{equation}
  |\theta_0|^2 + |\theta_1|^2 + |\theta_2|^2 + |\theta_3|^2 = 1.
  \label{eq:unitarity_full}
\end{equation}
A general element of $\mathbb{C}\otimes\mathbb{H}$ has $4 \times 2 = 8$ real degrees
of freedom (real and imaginary parts of each $\theta_\mu$).
Condition~\eqref{eq:unitarity_full} removes one real degree of freedom, leaving 7.
The group of such elements is $\mathrm{SU}(2) \times U(1)$, which has dimension 4;
modding out by the global $U(1)$ phase gives $\mathrm{SU}(2)$, dimension 3.

\subsection{Step 2: Restriction to $\mathrm{Im}(\mathbb{H})$}

On the imaginary part $\mathrm{Im}(\mathbb{H})$, we set $\theta_0 = 0$ and require
$\Theta = \theta_1\,\mathbf{i} + \theta_2\,\mathbf{j} + \theta_3\,\mathbf{k}$.
With real coefficients $\theta_a \in \mathbb{R}$ (the physical, non-complexified
sector of $\mathrm{Im}(\mathbb{H})$), unitarity requires
\begin{equation}
  \theta_1^2 + \theta_2^2 + \theta_3^2 = 1,
  \label{eq:unitarity_Im}
\end{equation}
i.e., $\Theta$ lies on the unit 2-sphere $S^2 \subset \mathrm{Im}(\mathbb{H}) \cong \mathbb{R}^3$.
The dimension of $S^2$ is 2, while $\dim_\mathbb{R}(\mathrm{Im}\,\mathbb{H}) = 3$.
The constraint removes one degree of freedom: $3 \to 2$.

\subsection{Step 3: Independent mode count after constraint}

The $N_{\mathrm{eff}} = 12$ modes arise from the product structure
$N_{\mathrm{eff}} = N_{\mathrm{phases}} \times N_{\mathrm{helicity}} \times N_{\mathrm{charge}} = 3 \times 2 \times 2$
(proved in \texttt{step3\_N\_eff\_result.tex}).
The constraint~\eqref{eq:unitarity_Im} acts on the $N_{\mathrm{phases}} = 3$
imaginary components.
After the constraint the independent phase modes span $S^2$, a 2-dimensional manifold.
The effective mode count becomes:
\begin{equation}
  N_{\mathrm{eff}}^{\mathrm{constrained}}
    = \frac{2}{3}\,N_{\mathrm{eff}}
    = \frac{2}{3} \times 12
    = 8.
  \label{eq:N_constrained}
\end{equation}
The corresponding $B$ coefficient would be (using the proved formula $B_0 = 2\pi N/3$):
\begin{equation}
  B_0(N=8) = \frac{2\pi \times 8}{3} = \frac{16\pi}{3} \approx 16.76.
  \label{eq:B_constrained}
\end{equation}
This value $\approx 16.76$ is much smaller than $B_{\mathrm{base}} = 41.57$.
The reduction factor is $2/3$, not $1/\sqrt{N_{\mathrm{eff}}}$.

\subsection{Step 4: Can any constraint give a $\sqrt{N_{\mathrm{eff}}}$ factor?}

The ratio $(3/2\pi)\sqrt{N_{\mathrm{eff}}} = 1.654$ requires that the constraint
multiplies $B_0$ by $\sqrt{N_{\mathrm{eff}}} \times (3/2\pi) \approx 1.654$.
This would correspond to the constrained system having an effective mode count of
\begin{equation}
  N_{\mathrm{eff}}^{\mathrm{target}} = N_{\mathrm{eff}}^{3/2} \cdot \frac{3}{2\pi N_{\mathrm{eff}}} \cdot \frac{3}{2\pi}
  = \frac{9 N_{\mathrm{eff}}^{1/2}}{4\pi^2}
  \approx \frac{9 \times 3.464}{39.48} \approx 0.790.
  \label{eq:N_target}
\end{equation}
A mode count less than 1 has no physical interpretation in a linear space.
The requirement that a unitarity constraint on a linear space of $N_{\mathrm{eff}}$
modes produces a ``number of modes'' less than 1 is impossible in the standard sense.

\begin{gap}[Approach D1 — Unitarity constraint]
\label{gap:D1}
  \textbf{[DEAD END].}
  The unitarity constraint $\Theta^\dagger\Theta = 1$ on $\mathrm{Im}(\mathbb{H})$
  reduces the independent real degrees of freedom from 3 to 2 (constraint on $S^2$),
  giving $N_{\mathrm{eff}}^{\mathrm{constrained}} = 8$ and $B \approx 16.76$.
  This is in the wrong direction: the constraint \emph{decreases} $B$ rather
  than increasing it toward $B_{\mathrm{base}} = 41.57$.
  Moreover, the target effective mode count $\approx 0.79 < 1$ is algebraically
  inconsistent with any linear constraint on a space of 12 modes.
  The analogy with $\mathrm{SU}(2)$ dimension reduction is valid but produces the
  wrong numerical factor.
\end{gap}

%--------------------------------------------------------------------
\section{Approach D2: Dimensional Transmutation on $\mathrm{Im}(\mathbb{H})$}
\label{sec:D2_dimtrans}
%--------------------------------------------------------------------

\textbf{Hypothesis.}
By analogy with QCD dimensional transmutation, an RG equation on
compact $\mathrm{Im}(\mathbb{H}) \cong \mathfrak{su}(2)$ flows from a UV scale
$\mu_{\mathrm{UV}}$ (set by $\psi$-compactification) to an IR scale
$\mu_{\mathrm{IR}}$ (lightest KK mode on $\mathrm{Im}(\mathbb{H})$).
The ratio $(\mu_{\mathrm{UV}}/\mu_{\mathrm{IR}})^\beta$ might equal
$\sqrt{N_{\mathrm{eff}}} \cdot (3/2\pi) = 1.6540$, requiring
$\beta \cdot \ln(\mu_{\mathrm{UV}}/\mu_{\mathrm{IR}}) = \ln(1.6540) = 0.5032$.

\subsection{Step 1: RG equation for $B(\mu)$ on $\mathrm{Im}(\mathbb{H})$}

The one-loop running of the $B$ coefficient under rescaling $\mu \to \mu e^t$ is
governed by the beta function of the quadratic coupling in $S[\Theta]$.
At one loop in the kinetic sector:
\begin{equation}
  \mu \frac{d B}{d\mu} = \beta_B \cdot B,
  \qquad
  \beta_B = \frac{N_{\mathrm{eff}}}{3} \cdot \frac{g^2}{(4\pi)^2},
  \label{eq:RG_B}
\end{equation}
where $g$ is the coupling constant of $S[\Theta]$ evaluated on $\mathrm{Im}(\mathbb{H})$.
Integrating from $\mu_{\mathrm{IR}}$ to $\mu_{\mathrm{UV}}$:
\begin{equation}
  B(\mu_{\mathrm{UV}}) = B(\mu_{\mathrm{IR}}) \cdot \exp\!\left(\beta_B
    \ln\frac{\mu_{\mathrm{UV}}}{\mu_{\mathrm{IR}}}\right)
  = B_0 \cdot \left(\frac{\mu_{\mathrm{UV}}}{\mu_{\mathrm{IR}}}\right)^{\!\beta_B}.
  \label{eq:B_running}
\end{equation}

\subsection{Step 2: Natural UV and IR scales from $\mathbb{C}\otimes\mathbb{H}$}

\textbf{UV scale.}
The $\psi$-compactification radius is $R_\psi$ (proved to set the UV cutoff).
The UV scale is therefore
\begin{equation}
  \mu_{\mathrm{UV}} = \frac{1}{R_\psi}.
  \label{eq:mu_UV}
\end{equation}

\textbf{IR scale.}
$\mathrm{Im}(\mathbb{H})$ is a 3-dimensional real vector space.
As a flat vector space $\mathbb{R}^3$, it has no compact geometry and no
natural IR scale independent of $R_\psi$.
When we exponentiate to obtain $\mathrm{SU}(2) \cong S^3$, the lightest mode
is the identity; the lightest \emph{non-trivial} mode has energy $1/R_{\mathfrak{su}(2)}$
where $R_{\mathfrak{su}(2)}$ is the curvature radius of $S^3$.

For the bi-invariant metric inherited from the quaternion inner product
$\langle \mathbf{i}, \mathbf{i}\rangle = 1$, the unit 3-sphere has
$R_{\mathfrak{su}(2)} = 1$ in natural quaternion units.
Thus:
\begin{equation}
  \mu_{\mathrm{IR}} = \frac{1}{R_{\mathfrak{su}(2)}} = 1.
  \label{eq:mu_IR}
\end{equation}

The ratio therefore satisfies:
\begin{equation}
  \frac{\mu_{\mathrm{UV}}}{\mu_{\mathrm{IR}}} = \frac{1}{R_\psi}.
  \label{eq:ratio_scales}
\end{equation}

\subsection{Step 3: Extracting $\beta_B$ from the condition}

The condition for dimensional transmutation to give the correct $B_{\mathrm{base}}$ is:
\begin{equation}
  \beta_B \cdot \ln\!\left(\frac{1}{R_\psi}\right) = \ln(1.6540) = 0.5032.
  \label{eq:beta_condition}
\end{equation}
This requires knowing $R_\psi$ in natural units.
In UBT, $R_\psi$ is set by calibration to $m_e$ (it is not derived from algebra
alone — see DERIVATION\_INDEX, row $R_\psi$).
Therefore $\ln(1/R_\psi)$ is a free parameter; no value of $\beta_B$ derived purely
from $S[\Theta]$ and $N_{\mathrm{eff}} = 12$ can simultaneously fix~\eqref{eq:beta_condition}
without using $R_\psi$ as calibrated input.

\subsection{Step 4: Intrinsic IR scale from algebra}

Could $\mathrm{Im}(\mathbb{H})$ itself supply an intrinsic IR scale
without reference to $R_\psi$?

As a Lie algebra, $\mathfrak{su}(2)$ carries the Killing form
$\kappa(X,Y) = -8\,\mathrm{Tr}(XY)$ for the standard normalisation
$\mathrm{Tr}(T_a T_b) = \delta_{ab}/2$.
The Killing form defines lengths but not a volume-scale independently
of the overall normalisation convention.
There is no canonical ``radius'' of $\mathfrak{su}(2)$ as an algebra;
only the associated \emph{group} $\mathrm{SU}(2) \cong S^3$ has a geometric radius,
and that radius depends on the metric normalisation choice (free parameter).

\begin{gap}[Approach D2 — Dimensional transmutation]
\label{gap:D2}
  \textbf{[DEAD END].}
  The RG flow of $B(\mu)$ on $\mathrm{Im}(\mathbb{H})$ requires two scales:
  $\mu_{\mathrm{UV}} = 1/R_\psi$ and $\mu_{\mathrm{IR}} = 1/R_{\mathfrak{su}(2)}$.
  \begin{enumerate}
    \item[(i)] $R_\psi$ is not algebraically fixed; it is currently calibrated to $m_e$
      and is therefore a free parameter.
    \item[(ii)] $R_{\mathfrak{su}(2)}$ depends on the overall normalisation of the
      $\mathrm{SU}(2)$ metric, which is itself a free parameter (only ratios within
      the algebra are fixed).
    \item[(iii)] Consequently, $\ln(\mu_{\mathrm{UV}}/\mu_{\mathrm{IR}})$ involves at
      least one free parameter, and the condition~\eqref{eq:beta_condition}
      can always be satisfied trivially by tuning.
      No parameter-free derivation of $B_{\mathrm{base}} = N_{\mathrm{eff}}^{3/2}$
      follows from this route.
  \end{enumerate}
  If an independent topological fixation of $R_\psi$ were found (see DERIVATION\_INDEX),
  then revisiting D2 would be meaningful, but the free-parameter obstruction would
  need to be resolved first.
\end{gap}

%--------------------------------------------------------------------
\section{Approach D3: Cartan--Killing Metric on $\mathfrak{su}(2)$}
\label{sec:D3_cartan_killing}
%--------------------------------------------------------------------

\textbf{Hypothesis.}
The mode-space $\mathrm{Im}(\mathbb{H}) \cong \mathfrak{su}(2)$ carries a natural metric
different from the standard Euclidean metric on $\mathbb{R}^3$: the Cartan--Killing
form $\kappa(X,Y) = \mathrm{Tr}(\mathrm{ad}_X \circ \mathrm{ad}_Y)$.
Using this metric in the vacuum polarisation integral $\Pi_{\mu\nu}$ might change the
$B$ coefficient from $B_0$ (Euclidean) to $B_{\mathrm{base}}$ (target).

\subsection{Step 1: Cartan--Killing form for $\mathfrak{su}(2)$}

The adjoint representation of $\mathfrak{su}(2)$ with generators
$T_a = -i\sigma_a/2$ ($\sigma_a$ = Pauli matrices) satisfies
$[T_a, T_b] = \varepsilon_{abc}\,T_c$.
The adjoint representation acts by $(\mathrm{ad}_{T_a})_{bc} = \varepsilon_{abc}$.

The Killing form is:
\begin{equation}
  \kappa_{ab}
    = \mathrm{Tr}(\mathrm{ad}_{T_a}\circ\mathrm{ad}_{T_b})
    = \sum_{c,d} (\varepsilon_{acd})(\varepsilon_{bcd})
    = \sum_{c,d} \varepsilon_{acd}\,\varepsilon_{bcd}.
  \label{eq:killing_form}
\end{equation}
Using the $\varepsilon$-contraction identity
$\sum_{c,d}\varepsilon_{acd}\,\varepsilon_{bcd} = 2\delta_{ab}$:
\begin{equation}
  \kappa_{ab} = 2\,\delta_{ab}.
  \label{eq:killing_result}
\end{equation}
(This matches the standard result $\kappa = -2\,\delta$ for the compact form with
the convention $[T_a,T_b]=\varepsilon_{abc}T_c$; the sign depends on convention.
For the purposes of metric determinant computation the sign is irrelevant.)

\subsection{Step 2: Effect on vacuum polarisation integral}

The one-loop vacuum polarisation in $d$ spatial dimensions with metric $g_{\mu\nu}$
on the mode space is
\begin{equation}
  \Pi \;\propto\; \int \frac{d^d k}{(2\pi)^d}\,
    \frac{\mathrm{Tr}[g^{-1}]}{(k^2 + m^2)^2}.
  \label{eq:Pi_general}
\end{equation}
The Killing form $\kappa_{ab} = 2\delta_{ab}$ differs from the Euclidean metric
$\delta_{ab}$ only by an overall factor of 2:
\begin{equation}
  \kappa^{-1}_{ab} = \tfrac{1}{2}\,\delta^{ab},
  \qquad
  \mathrm{Tr}[\kappa^{-1}] = \tfrac{d}{2} = \tfrac{3}{2}.
  \label{eq:killing_inverse}
\end{equation}

The contribution of the Killing metric relative to the Euclidean metric
therefore introduces a factor of $1/2$ in $\Pi$ and hence in $B$:
\begin{equation}
  B_{\mathrm{Killing}} = \frac{1}{2}\,B_0 = \frac{\pi N_{\mathrm{eff}}}{3}
    = \frac{\pi \times 12}{3} = 4\pi \approx 12.57.
  \label{eq:B_killing}
\end{equation}

This value $B_{\mathrm{Killing}} \approx 12.57$ is \emph{smaller} than $B_0 \approx 25.13$
and much smaller than $B_{\mathrm{base}} \approx 41.57$.

\subsection{Step 3: Mixed metric (Euclidean for momentum, Killing for mode structure)}

A physically more careful formulation considers the Euclidean metric for the
loop momentum $k$ (which propagates in the physical space $\mathrm{Im}(\mathbb{H})$)
and the Killing form only for the internal gauge-group factor in the trace.

For $\mathfrak{su}(2)$ with $N_c = 2$ colours, the standard one-loop gauge
contribution to $B$ from $N_{\mathrm{eff}}$ modes is
\begin{equation}
  B_{\mathrm{gauge}} = C_2(G)\,\frac{N_{\mathrm{eff}}}{3} \cdot (4\pi),
  \label{eq:B_gauge}
\end{equation}
where $C_2(G)$ is the quadratic Casimir of the adjoint representation.
For $\mathfrak{su}(2)$: $C_2(\mathrm{adj}) = 2$ (from $\mathrm{Tr}[T_a^{\mathrm{adj}}
T_b^{\mathrm{adj}}] = 2\delta_{ab}$ in the adjoint, consistent with
Eq.~\eqref{eq:killing_result}).

Substituting:
\begin{equation}
  B_{\mathrm{gauge}} = 2 \times \frac{12}{3} \times 4\pi = 8 \times 4\pi \approx 100.5.
  \label{eq:B_gauge_numerical}
\end{equation}
This is too large by a factor of $\approx 2.4$ compared to $B_{\mathrm{base}}$.

\subsection{Step 4: Check whether $\kappa_{ab}$ naturally replaces $\delta_{ab}$ in $B_0$}

The one-loop formula $B_0 = 2\pi N_{\mathrm{eff}}/3$ was derived using the
standard Euclidean metric $\delta_{ab}$ on $\mathrm{Im}(\mathbb{H}) \cong \mathbb{R}^3$.
Replacing $\delta_{ab} \to \kappa_{ab}/2$ (normalising Killing to unit at the root
length of $\mathfrak{su}(2)$) introduces a factor:
\begin{equation}
  \frac{\det(\kappa_{ab}/2)^{1/2}}{\det(\delta_{ab})^{1/2}}
  = \left(\frac{2}{2}\right)^{d/2}
  = 1^{3/2} = 1.
  \label{eq:det_ratio}
\end{equation}
No net factor: the normalised Killing form equals the Euclidean metric
when both are normalised so that the fundamental root has unit length.
The ratio $B_{\mathrm{Killing,norm}}/B_0 = 1$, which produces no correction at all.

\begin{gap}[Approach D3 — Cartan--Killing metric]
\label{gap:D3}
  \textbf{[DEAD END].}
  Three independent calculations show that the Cartan--Killing metric does not
  produce $B_{\mathrm{base}}$:
  \begin{enumerate}
    \item[(i)] Replacing $\delta_{ab}$ by $\kappa_{ab} = 2\delta_{ab}$ in $\Pi_{\mu\nu}$
      gives $B_{\mathrm{Killing}} = B_0/2 \approx 12.57$, moving in the
      \emph{wrong direction} (smaller, not larger).
    \item[(ii)] Using the $\mathfrak{su}(2)$ quadratic Casimir $C_2(\mathrm{adj})=2$
      in the standard gauge-group formula gives $B_{\mathrm{gauge}} \approx 100.5$,
      too large by a factor $\approx 2.4$.
    \item[(iii)] When both metrics are normalised to unit root length (canonical
      normalisation), the Killing form equals the Euclidean metric:
      $\kappa_{ab}^{\mathrm{norm}} = \delta_{ab}$, so the correction factor is exactly 1.
  \end{enumerate}
  In all three interpretations, the Cartan--Killing metric either decreases $B$
  or introduces a new free parameter (the normalisation of $\kappa$).
  No parameter-free derivation of $B_{\mathrm{base}} = N_{\mathrm{eff}}^{3/2}$ arises.
\end{gap}

%--------------------------------------------------------------------
\section{Conclusion: Results Table}
\label{sec:conclusion}
%--------------------------------------------------------------------

\begin{center}
\begin{tabular}{llll}
  \hline
  Approach & Result & Status & Direction \\
  \hline
  D1: Unitarity constraint on Im($\mathbb{H}$)
    & $B \to B_0 \times (2/3) \approx 16.8$
    & \textbf{[DEAD END]}
    & Wrong direction \\
  D2: Dimensional transmutation
    & $\beta_B \cdot \ln(\mu_{\mathrm{UV}}/\mu_{\mathrm{IR}}) = 0.503$?
    & \textbf{[DEAD END]}
    & Free parameter ($R_\psi$) \\
  D3: Cartan--Killing metric
    & $B_{\mathrm{Killing}} \approx 12.6$ or $100.5$
    & \textbf{[DEAD END]}
    & Wrong direction or too large \\
  \hline
\end{tabular}
\end{center}

\medskip\noindent\textbf{What was learned:}
\begin{itemize}
  \item \textbf{D1:} A unitarity constraint on $\mathrm{Im}(\mathbb{H})$
    decreases the effective mode count (from 12 to 8, not to $\sqrt{12} \approx 3.46$)
    and moves $B$ in the wrong direction.
    The target effective mode count $\approx 0.79$ implied by $B_{\mathrm{base}}$ is
    sub-unity and inconsistent with any real linear constraint.
  \item \textbf{D2:} Dimensional transmutation requires at least one free scale
    ($R_\psi$), so it cannot produce a parameter-free derivation of $B_{\mathrm{base}}$.
    If $R_\psi$ is fixed independently in the future, D2 should be revisited.
  \item \textbf{D3:} The Cartan--Killing metric on $\mathfrak{su}(2)$ is proportional
    to the Euclidean metric; its normalised form gives no correction.
    Unnormalised choices either halve or multiply $B$, neither matching $B_{\mathrm{base}}$.
  \item \textbf{Overall:} The gap $B_0 \to B_{\mathrm{base}}$ (equivalently,
    $\Delta d = 0.405$) does not arise from standard perturbative or semi-perturbative
    mechanisms derived from $N_{\mathrm{eff}} = 12$ and $\dim_\mathbb{R}(\mathrm{Im}\,\mathbb{H}) = 3$
    alone. A genuinely new structural input appears to be required.
\end{itemize}

\textbf{Remaining open directions} (not yet explored):
\begin{itemize}
  \item Instanton contributions on $\mathrm{SU}(2)$
    (exponentially small in weak coupling, but non-perturbative).
  \item Index theorem or anomaly constraints on the mode space.
  \item Holomorphic factorisation in $\mathbb{C}\otimes\mathbb{H}$ reducing the
    \emph{complex} mode count.
  \item Non-commutative geometry approach to the spectral triple of $\mathbb{C}\otimes\mathbb{H}$.
\end{itemize}

%--------------------------------------------------------------------
\section{In Plain Language}
\label{sec:plain}
%--------------------------------------------------------------------

\textbf{In plain language:}
Three non-perturbative routes were explored to explain why the $B$ coefficient
jumps from the proved value $B_0 \approx 25.13$ to the conjectured value
$B_{\mathrm{base}} \approx 41.57$.
The unitarity-constraint approach (D1) was ruled out because imposing a constraint
on the biquaternion modes actually makes $B$ smaller, not larger.
The dimensional-transmutation approach (D2), inspired by QCD, was blocked by the
absence of an algebraically-fixed UV/IR scale ratio — the $\psi$-radius $R_\psi$
remains a calibrated rather than derived parameter.
The Cartan--Killing-metric approach (D3) was ruled out because the natural
$\mathfrak{su}(2)$ metric is proportional to the ordinary Euclidean metric and
produces either a decrease or an overshoot in $B$, depending on normalisation.
The gap $B_0 \to B_{\mathrm{base}}$ therefore remains open after eight independent
dead ends; any successful approach will likely require a structural input beyond
mode counting and standard perturbative geometry.

\ifdefined\INCLUDEMODE\else
\end{document}
\fi
